\chapter{Conclusion générale}

\begin{quote}
    \textit{On peut, bien entendu, comprendre de la même façon le développement de l'art, et de la science, de la religion et de toutes formes de culture, élevées ou inférieures. En tirant une leçon des efforts des uns et des autres et en appréciant leurs nombreuses contributions, les êtres humains construisent progressivement des systèmes de connaissances et de croyances ; ils élaborent des techniques reconnues et des styles élaborés de sensibilité et d'expression. Dans ces exemples, le but commun est souvent profond et complexe, étant défini par la tradition correspondante, artistique, scientifique ou religieuse et, pour le comprendre, il faut souvent des années d'étude et de discipline. L'essentiel est de partager  une fin et d'être d'accord sur les moyens de l'atteindre afin de permettre la reconnaissance publique des réalisations de chacun. Quand cette fin est atteinte, tous trouvent de la satisfaction dans la même chose ; et ce fait ainsi que la complémentarité du bien des individus renforcent les liens de la communauté.}

    John Rawls, \textit{Théorie de la justice}, chapitre 79 \og L'idée d'union sociale \fg, page 569, éditions du Point, 1971.
    
\end{quote}

Voici venu l'instant de clore ce travail sur l'univers foisonnant de la culture du piratage culturel, des individus qui vivent cette culture au jour le jour, des communautés qu'ils forment, et des actions que ces communautés entreprennent pour faire vivre l'idée que tout ce qui a été digne d'être créé mérite d'être conservé et accessible à tous. 

Dans notre introduction, nous avons défini le sujet de notre étude comme étant les communautés patrimoniales pirates : une organisation sociale existant par et pour une relation avec du patrimoine immatériel. Nous avons décidé d'étudier les expériences, les émotions et les actions des individus qui ont comme point commun l'utilisation et l'attachement à des services de bibliothèques et d'archives pirates. Nous avons également voulu étudier la conséquence du tournant numérique sur la question de la reproductibilité et de la propriété intellectuelle quand on la sort du contexte institutionnel et entrepreneurial. 

Nous nous sommes questionnés ainsi : Comment les dynamiques culturelles et les organisations sociales permettent-elles l'émergence d'une pratique archivistique populaire, autogestionnaire et clandestine ? Nous avons exploré la manière dont les individus découvrent et intègrent la culture et le fonctionnement des dépôts pirates. Nous avons ensuite cherché le point de rencontre entre l'expérience des individus et l'organisation de guérilla des bibliothèques pirates. Enfin, dans notre troisième chapitre, nous avons réalisé une typologie des différentes traductions des communautés patrimoniales pirates, du dépôt personnel aux bibliothèques d'Alexandrie clandestines contemporaines.\\

Une partie de la réponse se trouve au cœur de la nouvelle nature qu'ont acquise les objets culturels, avec la conséquence sur nos horizons mentaux de la massification de la capacité de copie des individus. Cette transformation de la capacité et du coût de copie entre en contradiction avec le mode de distribution et le fonctionnement économique de nos sociétés. Les pratiques de consommations culturelles actuelles sont perçues comme aliénantes et injustes. Où la culture et le divertissement, par extension sa patrimonialisation, ont été confisqués par un système marchand constitué du capitalisme centré autour de deux mannes de revenus principales : la surveillance et le service \footnote{Soshana Zuboff, \textit{L'âge du capitalisme de surveillance,}, éditions Zulma, 2019, \og Nouvel ordre économique qui revendique l'expérience humaine comme matière première gratuite à des fins de pratiques commerciales dissimulées d'extraction, de prédiction et de vente \fg }, les deux demandant un fonctionnement cloisonné, sécurisé. Cette mécanique est alors vécue comme une dépossession, qui pousse les individus à mettre en place un espace en dehors des codes sociaux classiques, un espace d'échange de l'information qui ne se limite qu'à la bande passante et à la capacité de stockage. Nous l'avons étudié chez nos enquêtés, tous trouvent un mécontentement, une injustice dans la manière dont le système culturel fonctionne, même s'ils n'y trouvent pas de considérations politiques, ce qui justifie leur passage à l'action en tant que pirates. Cette émotion, que l'on retrouve dans les écrits de D.Fabre, E. Bermès ou A. Kosnik. Ces dépôts pirates sont chargés d’autant de projets utopiques que d'espaces inquiétants et dangereux. Ces dépôts pirates sont chargés d’autant de projets utopiques que d'espaces inquiétants et dangereux. Le piratage offre un espace alternatif, qui, en plus de l’appât du gain économique, permet de toucher du doigt une forme de distribution alternative. Pour se retrouver dans les méandres des forums et des bibliothèques, il faut alors apprendre à maîtriser la culture, les interactions sociales, les nouveaux codes, les outils et l'esthétique particulière générée par la clandestinité et les représentations qu'ont les communautés d'elles-mêmes. Cet apprentissage est un phénomène social qui peut connecter des individus dans le monde physique (socialisation par les pairs, figure tutélaire), ou connecter l'individu à des ressources et des communautés qui permettent un apprentissage en autonomie. Une fois les règles et les outils maîtrisés, l'individu va mettre en place une routine et des habitudes de consommation incluant le piratage. Une rupture avec la pratique peut se produire si ces habitudes sont désorganisées (fermeture d'un site très utilisé, mise hors ligne d'un outil, moins de temps à dédier au piratage à cause d'un emploi, une situation économique qui permet de sortir du circuit de piratage). Une désorganisation des habitudes peut également pousser l'individu à évoluer et acquérir de nouvelles compétences. Ces compétences peuvent être également utiles et devenir une source de fierté dans d'autres interactions que la sphère culturelle pirate. L'individu peut alors devenir lui-même une passerelle pour ses pairs. \\

Deuxièmement, les bibliothèques pirates les plus connues comme Sci-Hub ou LibGen sont les éléments les plus visibles et médiatisés de ce que nous avons nommé une culture d'archivistique émergente. Cette culture tire son origine du lien émotionnel que les individus investissent dans le patrimoine culturel numérique. Ces dépôts sont à la fois le symbole de l'existence d'une fragilité concernant la conception de la propriété dans nos sociétés, et la tentative de créer une alternative, dans un espace non modéré. Cette dynamique de distribution culturelle parallèle a un impact extrêmement fort sur l'accès à la culture de toute une partie de la population la plus précaire, mais également des ramifications très profondes sur des pratiques professionnelles où des secteurs économiques entiers, comme par exemple le monde du développement de modèles d'intelligences artificielles génératives. En conclusion, l'étude du phénomène des bibliothèques pirates et des communautés patrimoniales pirates met en lumière les tensions inhérentes entre les modèles économiques actuels et la diffusion alternative de la culture, ainsi que les nouvelles compétences, solidarités et dangers qui en émergent.\\

De nombreux points traités dans ce travail pourraient être approfondis avec plus de temps. Nous n'avons pu nous entretenir qu'avec cinq enquêtés, et une grande partie de nos sources ne sont que quelques exemples perdus dans l'infinité du Web. Également, de nombreux projets d'exploitations n'ont pas été incorporés à notre travail. Nous avons également fait le choix de nous limiter au versant le plus proche de nos situations sociales, mais il existe toute une littérature sur les cultures pirates en dehors du nord économique. Nous n'avons pu également qu'effleurer certaines communautés comme les phénomènes low tech, les communautés DIY/Hacktivistes et les acteurs du \textit{Right to repair} et \textit{Right to Hack}. Ce travail ne peut pas proposer de réponses, seulement de nouvelles pistes à explorer. Pour prolonger les questionnements de ce mémoire après sa lecture, notre souhait est l'ouverture et l'approfondissement de la compréhension des pratiques auto-organisées et de leur influence sur l'accès à la culture, la conservation du patrimoine, mais également dans tous les domaines où les règles peuvent être tordues et réinventées pour des causes justes.  
